% Image credits for the photographs and AI illustrations of Book 4
% (University Chemistry, Year 3). Machine-readable license records live in
% images/book4/CREDITS.md; keep the two files in sync. The Book 4 writer
% owns this file: add one \item per photograph (in an itemize, after the
% first paragraph) as photographs land.
\chapter*{Image Credits}

The drawings in this book are original. Photographs, where used, are used
under their respective free licenses, with thanks to their authors; full
source links and license URLs for every photograph are listed in
\texttt{images/book4/CREDITS.md} in the book's repository.
The hazard pictograms are those of the Globally Harmonized System of
Classification and Labelling of Chemicals, as drawn for the
\texttt{ghsystem} package of \TeX\ Live.

\bigskip
Alongside the photographs and the original drawings, some figures are
AI-generated illustrations, produced with OpenAI image generation from
prompts written by the book's editors, and reviewed for chemical
accuracy before inclusion. The full prompt for every generated image is
recorded in \texttt{images/book4/ai/PROMPTS.md} in the repository.
\bigskip
\noindent Photographs, by chapter:
\begin{itemize}\small
  \item Ch.~1: Erwin Schr\"odinger, 1933, Nobel Foundation, public domain
        (Wikimedia Commons).
  \item Ch.~4: snow crystals, Wilson A. Bentley (about 1902), public domain
        (Wikimedia Commons).
  \item Ch.~6: radio antennas on a high plateau, ESO/B.~Tafreshi
        (twanight.org), CC BY 4.0 (Wikimedia Commons).
  \item Ch.~7: fluorite in daylight and under ultraviolet light, by Masha
        Milshina, CC BY 4.0 (Wikimedia Commons).
  \item Ch.~8: an NMR laboratory, by Shandchem, CC BY 2.0 (Wikimedia
        Commons).
  \item Ch.~9: William Lawrence Bragg, Nobel Foundation, public domain
        (Wikimedia Commons).
  \item Ch.~10: the grave of Ludwig Boltzmann, by Feldkurat Katz,
        CC BY-SA 4.0 (Wikimedia Commons).
  \item Ch.~13: Belousov--Zhabotinsky spirals, by Stephen Morris, CC BY 2.0
        (Wikimedia Commons).
  \item Ch.~14: the Antarctic ozone hole, NASA, public domain (Wikimedia
        Commons).
  \item Ch.~16: electron micrograph of a catalytic-converter honeycomb, by
        Galadrid, CC BY-SA 4.0 (Wikimedia Commons).
  \item Ch.~17: sunbeams and the Tyndall effect, by Kevin c higgins, CC BY-SA 4.0
        (Wikimedia Commons).
  \item Ch.~18: ruby and emerald crystals, by Robert M. Lavinsky, CC BY-SA 3.0
        (Wikimedia Commons).
  \item Ch.~20: crystals of ferrocene, by Andrea Sella, CC BY 4.0 (Wikimedia
        Commons).
  \item Ch.~22: a single-crystal silicon boule, by Sebastian Wallroth, public
        domain (Wikimedia Commons).
  \item Ch.~23: a flower on silica aerogel, NASA/JPL, public domain; the
        Lycurgus cup, by Marie-Lan Nguyen, CC BY 2.5 (Wikimedia Commons).
  \item Ch.~24: an Atlantic horseshoe crab, by Kaldari, CC0 (Wikimedia
        Commons).
  \item Ch.~27: pyrethrum daisies, by Soramimi, CC BY-SA 4.0 (Wikimedia
        Commons).
  \item Ch.~30: bark of a Pacific yew, by JOE BLOWE (theforestprimeval),
        CC BY-SA 2.0 (Wikimedia Commons).
  \item Ch.~32: an algal bloom seen by the Landsat 8 satellite, USGS and NASA,
        public domain.
  \item Ch.~33: a double vacuum and inert-gas manifold, by Polimerek, CC BY-SA
        3.0 (Wikimedia Commons).
\end{itemize}
\clearpage
